\documentclass[11pt]{article}

\usepackage[T1]{fontenc}
\usepackage[utf8]{inputenc}
\usepackage{amsmath,amssymb,mathtools}
\usepackage{booktabs}
\usepackage{geometry}
\usepackage{graphicx}
\usepackage{hyperref}
\usepackage{microtype}
\usepackage{xcolor}

\geometry{margin=27mm}
\setlength{\emergencystretch}{1em}
\hypersetup{colorlinks=true,linkcolor=blue!55!black,urlcolor=blue!55!black}
\newcommand{\R}{\mathbb{R}}
\newcommand{\E}{\mathbb{E}}
\newcommand{\Cov}{\operatorname{Cov}}
\newcommand{\diag}{\operatorname{diag}}
\newcommand{\bits}{\text{ bits}}

\title{Latent Reconciliation Histories, Missing Information,\\
and Coordinate Changes in GPURec}
\author{A mathematical synthesis of two optimisation-agent transcripts}
\date{6 September 2026}

\begin{document}
\maketitle

\begin{abstract}
This report explains what two agents learned about gene-wise optimisation in GPURec's
duplication--transfer--loss reconciliation model.  The three base-2 log rates are already the
natural parameters of the complete-data model, whose negative log-likelihood is convex.  The
observed objective can nevertheless be non-convex because reconciliation histories are latent:
its curvature is complete-data multinomial information minus a posterior covariance of event
counts, with an additional contribution from conditioning on family survival.  Nonlinear
coordinates such as square-root rates substantially improve the measured local geometry, but
they do not reduce the production cost when curvature is carried by BFGS updates.  The measured
cost remains about eleven clade-weighted population passes.  The EM-focused agent derived an EM
iteration from the same structure, but its transcript ended before count extraction,
missing-information fractions, or EM convergence were measured; those quantities therefore
remain open rather than being inferred from the coordinate experiment.
\end{abstract}

\section{Model and notation}

For one gene family, let
\[
  \theta=(\theta_D,\theta_L,\theta_T)
  = (\log_2 D,\log_2 L,\log_2 T)
\]
denote the duplication, loss, and transfer log rates.  Speciation is the reference category with
fixed logit zero.  With
\[
  Z(\theta)=1+2^{\theta_D}+2^{\theta_L}+2^{\theta_T},
\]
the four event probabilities are
\begin{equation}
  p_S=\frac1Z,
  \qquad
  p_k=\frac{2^{\theta_k}}Z,
  \quad k\in\{D,L,T\}.
  \label{eq:softmax}
\end{equation}
The transfer probability is subsequently distributed among permissible recipient species by a
factor that does not depend on the three rates.  The fitted rates are constrained to
$10^{-6}\leq D,L,T\leq 2$, or approximately
$-19.93\leq\theta_k\leq1$.

A \emph{complete history} $h$ records every event of every lineage, including extinct lineages.
Write $n_k(h)$ for its number of events of type $k$ and
$n_\bullet(h)=\sum_{k\in\{S,D,L,T\}}n_k(h)$.  Rate-independent factors, including gene-tree split
weights and recipient choices, are collected into $w_h$.  Thus
\begin{equation}
  P(h\mid\theta)
  =w_h\prod_{k\in\{S,D,L,T\}}p_k(\theta)^{n_k(h)}.
  \label{eq:history-probability}
\end{equation}

For observed gene-tree data $y$, let $\mathcal H_y$ be the compatible histories and define the
unconditioned numerator
\[
  A_y(\theta)=\sum_{h\in\mathcal H_y}P(h\mid\theta).
\]
If $q(\theta)=P_\theta(\text{family survives})$, GPURec uses the conditional likelihood
\begin{equation}
  L_y(\theta)=\frac{A_y(\theta)}{q(\theta)},
  \qquad
  \ell(\theta)=-\log_2 A_y(\theta)+\log_2 q(\theta).
  \label{eq:conditioned-likelihood}
\end{equation}
The survival probability is obtained from extinction probabilities computed by a fixed point.

\section{Why the log rates are the natural coordinates}

For a fixed complete history, substitution of \eqref{eq:softmax} into
\eqref{eq:history-probability} gives
\begin{align}
  -\log_2 P(h\mid\theta)
  &=n_\bullet(h)\log_2 Z(\theta)
    -\sum_{k\in\{D,L,T\}}n_k(h)\theta_k
    -\log_2 w_h.                                      \label{eq:complete-nll}
\end{align}
This is a log-sum-exp minus a linear term.  Consequently the complete-data negative
log-likelihood is convex.  More precisely, for
\[
  p=(p_D,p_L,p_T)^\mathsf T,
  \qquad
  C(p)=\diag(p)-pp^\mathsf T,
\]
its Hessian in base-2 log-rate coordinates is
\begin{equation}
  \nabla_\theta^2[-\log_2P(h\mid\theta)]
  =n_\bullet(h)\ln(2)\,C(p)\succeq0.                 \label{eq:complete-hessian}
\end{equation}
The usual natural parameters are $\eta_k=\ln(D_k)=\ln(2)\theta_k$; multiplication by the common
positive constant $\ln(2)$ is immaterial.  Hence the existing log rates are already the natural
coordinates, and the event counts are sufficient statistics.  A nonlinear power transform may
be computationally useful, but it is not more intrinsic to the complete-data statistical model.

\section{Where observed-data non-convexity comes from}

Marginalisation over histories changes the curvature.  Define the complete-history score statistic
\[
  s(h)=
  \begin{pmatrix}n_D(h)\\n_L(h)\\n_T(h)\end{pmatrix}
  -n_\bullet(h)p.
\]
For the unconditioned data term, the missing-information identity gives
\begin{equation}
  H_y
  =\ln(2)\left\{
     \E_y[n_\bullet]C(p)-\Cov_y[s(h)]
   \right\}.                                         \label{eq:louis-data}
\end{equation}
Here the expectation and covariance are under the posterior distribution of histories conditional
on $y$ at the current parameter.  The first term is the expected complete-data information.  It is
positive semidefinite and depends only on one expected total count and the four event probabilities.
The second term is the information missing because the history was not observed.  It is a posterior
covariance and is also positive semidefinite, but it is subtracted.  If alternative histories have
very different event counts, this covariance can dominate the multinomial term and make $H_y$
indefinite.

Survival conditioning in \eqref{eq:conditioned-likelihood} must also be included.  One direct form is
\begin{equation}
  H
  =\ln(2)\left[
    \bigl\{\E_y[n_\bullet]C-\Cov_y(s)\bigr\}
    -\bigl\{\E_{\rm surv}[n_\bullet]C-\Cov_{\rm surv}(s)\bigr\}
  \right].                                           \label{eq:conditioned-information}
\end{equation}
Thus the survival histories enter with the opposite sign.  The EM transcript also identified a
useful equivalent augmentation.  If $e=1-q$ is the extinction probability, then
\[
  \frac1q=\frac1{1-e}=\sum_{m=0}^{\infty}e^m.
\]
The conditional likelihood may therefore be augmented by an observed family together with a
geometrically distributed number of unobserved extinct ``ghost'' families.  In this augmentation
all expected event counts are positive and an identity of the form
$H=I_c-I_m$ again holds with $I_c\succeq0$ and $I_m\succeq0$.  The direct survival-difference
counts and ghost-augmented counts give the same score: their discrepancy is proportional to the
probability vector and is annihilated by the softmax Jacobian.  The decomposition into complete and
missing information depends on the chosen augmentation, whereas the observed Hessian $H$ does not.

\subsection{Why a coordinate map cannot remove the statistical ambiguity}

Let $\phi$ be any smooth invertible coordinate system, write $\theta=t(\phi)$, let
$J=\partial\theta/\partial\phi$, and put $g=\nabla_\theta\ell$.  The transformed Hessian is
\begin{equation}
  \nabla_\phi^2\ell
  =J^\mathsf T HJ+sum_{k\in\{D,L,T\}}g_k\nabla_\phi^2t_k.
  \label{eq:hessian-transform}
\end{equation}
At a stationary point $g=0$, only the congruence $J^\mathsf THJ$ remains.  Sylvester's law of
inertia then preserves the numbers of positive and negative eigenvalues.  Away from a stationary
point, a nonlinear map can add helpful curvature through the second term in
\eqref{eq:hessian-transform}; it cannot eliminate the posterior competition itself, and that added
curvature vanishes at the optimum.

For example, power coordinates $\phi_k=r_k^\alpha=2^{\alpha\theta_k}$ have inverse
\[
  \theta_k=\frac{\ln\phi_k}{\alpha\ln2},\qquad
  \frac{d\theta_k}{d\phi_k}=\frac1{\alpha\ln2\,\phi_k},\qquad
  \frac{d^2\theta_k}{d\phi_k^2}=-\frac1{\alpha\ln2\,\phi_k^2}.
\]
This explains both observations in the experiment: square-root rates can add useful curvature far
from the optimum, while rate coordinates become numerically delicate near the $10^{-6}$ floor
because the second derivative of the inverse map is very large.

\section{The principled latent-variable iteration}

The expectation--maximisation (EM) algorithm follows directly from
\eqref{eq:complete-nll}.  Its E-step computes posterior expected event counts
\[
  N_k=\E[n_k\mid y,\theta^{(r)}],
  \qquad N_\bullet=\sum_kN_k.
\]
Using the positive ghost-family augmentation for survival conditioning, the M-step minimises the
expected complete-data NLL and has the multinomial closed form
\begin{equation}
  p_k^{(r+1)}=\frac{N_k}{N_\bullet},
  \qquad
  \theta_k^{(r+1)}=\log_2\frac{N_k}{N_S},
  \quad k\in\{D,L,T\}.                               \label{eq:em-update}
\end{equation}
If instead one uses data-side minus survival-side counts directly, the same formula applies when
all effective counts are positive.  A nonpositive effective count sends the unconstrained optimum
to a boundary; an active-set M-step or a Newton minimisation of the convex surrogate is then the
appropriate implementation.

The box-constrained M-step also remains analytic.  Suppose a set $B$ of rates is pinned at log-rate
bounds $b_k$, and let $F$ denote the free rates.  With
\[
  c=1+\sum_{k\in B}2^{b_k},
\]
the free solution is
\begin{equation}
  \theta_k^{(r+1)}
  =\log_2\left(
       \frac{N_kc}{N_S+\sum_{j\in B}N_j}
    \right),\qquad k\in F,                          \label{eq:box-em}
\end{equation}
followed by the usual active-set and Karush--Kuhn--Tucker sign checks.

The existing reverse pass already accumulates adjoints with respect to the four free event
log-probabilities before folding them through the softmax.  If
$a_k=\partial\ell/\partial\log_2p_k=-N_k$, then the ordinary log-rate gradient is
\begin{equation}
  g_k=a_k-p_k\sum_j a_j=p_kN_\bullet-N_k.             \label{eq:gradient-counts}
\end{equation}
Thus the E-step should require essentially the same backward computation as a gradient pass.

Near an optimum, define the complete information $I_c$, observed information $I_o=H$, and missing
information $I_m=I_c-I_o$.  The local EM error-propagation matrix is
\begin{equation}
  R_{\rm EM}=I_c^{-1}I_m,                             \label{eq:em-rate}
\end{equation}
up to a similar symmetric representation.  Its spectral radius is the fraction of missing
information and determines the linear convergence rate.  A generalized eigenvalue larger than
one is equivalent to negative observed curvature in the corresponding direction.  This supplies
a statistical explanation for why BFGS--Newton can be viewed as an acceleration: it attempts to
learn $I_o=I_c-I_m$, whereas plain EM uses the inexpensive $I_c$.

\paragraph{Evidence boundary.}
The EM transcript stopped immediately after writing this derivation.  It did not create the planned
count-extraction script, validate \eqref{eq:gradient-counts}, measure the missing-information
fractions, or run EM, SQUAREM, or a hybrid method.  The report therefore makes no numerical claim
about EM's actual convergence on the Coleman data.

\section{What the coordinate experiment measured}

The completed experiment used the first 500 Coleman gene families and exact $3\times3$ Hessians at
the common start, after three Adam steps, and at the fitted optimum.  A distance is the largest of
the three absolute log-rate gaps from the optimum, followed by the median across families.  The
post-Adam median distance was $2.48$ log-rate units and the 90th percentile journey was about $8.5$.

\subsection{Direct evidence for competing reconciliation explanations}

At the post-Adam point, duplication was above its eventual optimum in $72\%$ of families and
transfer was below its optimum in $93\%$.  The median transfer gap was $2.09$ log-rate units, or
approximately a fourfold rate increase.  Nevertheless, the local duplication gradient requested
an \emph{increase} in $69\%$ of families.  At artificially low transfer, duplication temporarily
explains discordant gene histories; as transfer rises, those duplications are no longer needed.

This competition also appears in exact negative-curvature eigenvectors.  At the post-Adam point,
$320/500=64\%$ of Hessians were indefinite.  The most negative direction had mean squared weights
$83\%$ duplication, $16\%$ transfer, and $1\%$ loss; its duplication and transfer components had
opposite signs in $95\%$ of indefinite families.  At the common start, $410/500=82\%$ were
indefinite and the corresponding direction was $63\%$ loss, $31\%$ transfer, and $6\%$
duplication.  These measurements are exactly the pattern predicted by the covariance term in
\eqref{eq:louis-data}: histories exchange duplication or loss events for transfer events.

\subsection{Coordinate screening and exact-Hessian trajectories}

Table~\ref{tab:coordinates} reports the principal coordinate results.  ``Indefinite'' gives the
percentage at the start, post-Adam point, and optimum.  ``Drift'' is the median relative Frobenius
distance between the post-Adam and optimum Hessians in the stated coordinates.  Multi-step entries
are GPU-evaluated, accepted exact-Hessian steps; one-step-only entries are from the CPU screening.

\begin{table}[htbp]
\centering
\small
\setlength{\tabcolsep}{4.5pt}
\caption{Geometry and median distance to the optimum for candidate coordinates.}
\label{tab:coordinates}
\resizebox{\textwidth}{!}{%
\begin{tabular}{lccrrrr}
\toprule
Coordinates & Indefinite start/Adam/opt. & Drift & Step 1 & Step 2 & Step 3 &
  NLL excess after step 3 \\
\midrule
Log$_2$ rates, production radius rule & $82/64/0.2\%$ & $1.01$ & $1.65$ & $1.40$ & $0.75$ & $697\bits$ \\
Log$_2$ rates, radius fixed at 2       & $82/64/0.2\%$ & $1.01$ & $1.65$ & $0.98$ & $0.47$ & $876\bits$ \\
Square-root rates                      & $26/10/0\%$    & $0.90$ & $1.27$ & $0.55$ & $0.07$ & $141\bits$ \\
Cube-root rates                        & $47/23/0.2\%$  & $0.56$ & $1.30$ & $0.61$ & $0.08$ & $46\bits$ \\
Square-root event probabilities        & $47/10/0.2\%$  & $0.86$ & $1.40$ & $0.49$ & $0.04$ & $55\bits$ \\
Linear rates                           & $35/33/36\%$   & $3.70$ & $1.73$ & --- & --- & --- \\
Event-probability simplex              & $16/12/11\%$   & $3.02$ & $1.77$ & --- & --- & --- \\
Per-event log odds                     & $86/53/0.2\%$  & $0.94$ & $1.70$ & --- & --- & --- \\
Log total rate plus event shares       & $30/7/6\%$     & $4.70$ & $1.98$ & --- & --- & --- \\
Oracle whitening by optimum Hessian    & $82/64/0.2\%$ & $2.33$ & $2.57$ & --- & --- & --- \\
\bottomrule
\end{tabular}}
\end{table}

Square-root rates reduced post-Adam indefiniteness from $64\%$ to $10\%$ and reduced the median
optimum condition number from about $80$ to $5.4$; square-root probabilities achieved a condition
number near $6.1$.  Cube-root rates gave the lowest Hessian drift, $0.56$, and the smallest measured
quadratic-model error in the screening.  Exact-Hessian trajectories consequently improved from
\[
  2.48\to1.65\to0.98\to0.47
  \quad\text{(fixed-radius log rates)}
\]
to
\[
  2.48\to1.27\to0.55\to0.07
  \quad\text{(square-root rates)}.
\]
The apparent indefiniteness of linear-rate and raw-probability coordinates at the constrained
optimum is not evidence of an interior saddle.  Many points lie at a rate bound, so the residual
gradient term in \eqref{eq:hessian-transform} need not vanish; the inverse map also has derivatives
of order $r^{-2}$ near the floor.

The production trust radius itself imposes a lower bound on one-step progress.  A perfect
straight-line displacement toward the optimum, capped at two log-rate units, would still leave a
median distance of $0.64$.  Moreover, after three exact steps only $5$--$11\%$ of families in any
coordinate system passed the actual certificate
$\lVert g_{\rm projected}\rVert_\infty<10^{-3}$.  Median geometric proximity is therefore not the
same as convergence of the costly tail, particularly because large-clade families converge last.

\subsection{Why improved geometry did not improve production cost}

An exact Hessian costs about $7.7$ whole-population gradient passes.  Three exact-Hessian steps in
square-root coordinates therefore cost approximately
$3(7.7+1)\simeq26$ passes.  Taking one exact Hessian and then carrying it with BFGS cost $15.8$
clade-weighted passes, compared with about $11.0$ for the production replay.

The decisive production-like test used one gradient pass per round and a BFGS-carried curvature.
Table~\ref{tab:bfgs} shows the state after twelve rounds.  The native seeds were constructed in
each coordinate system from the same start-to-post-Adam gradient pair.

\begin{table}[htbp]
\centering
\small
\caption{Production-like BFGS replay after twelve rounds.}
\label{tab:bfgs}
\begin{tabular}{lrrrr}
\toprule
Run & Frozen families & Live clades & Weighted cost & NLL excess \\
\midrule
Log$_2$ rates, production seed & $63\%$ & $21\%$ & $11.0$ & $1.7\bits$ \\
Log$_2$ rates, native seed     & $58\%$ & $32\%$ & $11.01$ & $2.0\bits$ \\
Square-root rates, native seed & $37\%$ & $73\%$ & $11.76$ & $22.7\bits$ \\
Cube-root rates, native seed   & $73\%$ & $24\%$ & $11.25$ & $0.9\bits$ \\
Square-root probabilities, native seed & $80\%$ & $19\%$ & $11.10$ & $34.3\bits$ \\
Square-root rates, exact seed & $80\%$ & $6\%$ & $15.8$ & --- \\
\bottomrule
\end{tabular}
\end{table}

Thus the useful nonlinear maps changed curvature but not end-to-end cost under the curvature scheme
that production can afford.  A rank-two BFGS update still has to learn a drifting $3\times3$
matrix from successive gradient pairs; the trust radius limits long moves; and the clade-weighted
tail dominates wall time.  The non-convexity is also structural rather than a mere scaling defect:
the posterior over histories changes as transfer rises, so the Hessian changes with it.

\section{Conclusions}

The combined mathematical and empirical conclusions are:
\begin{enumerate}
  \item The base-2 log rates are the model's natural coordinates and make every fixed-history
        complete-data objective convex.
  \item Marginalising reconciliation histories subtracts a posterior score covariance from the
        multinomial information.  Survival conditioning contributes an analogous term.  This
        missing information is the source of the measured duplication--transfer and
        loss--transfer negative directions.
  \item Nonlinear power coordinates can improve off-optimum curvature through a
        gradient-dependent Hessian term.  Square-root and cube-root variants markedly improve
        exact-Newton trajectories, but no smooth map removes the latent ambiguity, and at a true
        stationary point coordinate changes preserve inertia.
  \item With production's BFGS-carried curvature, the tested maps cost about
        $11.0$--$11.8$ clade-weighted passes, while an exact-Hessian seed is too expensive.
        Reparameterisation therefore does not cut the Newton phase from roughly eleven passes to a
        few.
  \item EM is the principled remaining experiment: its M-step is analytic and its E-step should be
        obtainable from the existing adjoints, but neither its counts nor its missing-information
        fraction nor its convergence trace was measured in the supplied transcript.
\end{enumerate}

\paragraph{One-sentence verdict.}
The agents learned that the optimisation difficulty is caused primarily by posterior uncertainty
between competing reconciliation histories, not by a poor choice of rate coordinates: nonlinear
coordinates improve exact local geometry, but they do not reduce affordable production passes,
and the proposed EM alternative remains theoretically motivated but experimentally untested.

\section*{Provenance}

This synthesis is based on the two supplied JSONL task outputs.  The completed reparameterisation
study reported that it changed no repository source files and stored its scripts and measurements
under the scratchpad's \texttt{reparam/} directory.  The EM task wrote only its derivation record
under \texttt{em/PROGRESS.md} before its transcript ended.  This document is the only new source
artifact created for the present request.

\begin{thebibliography}{9}
\bibitem{DLR1977}
A. P. Dempster, N. M. Laird, and D. B. Rubin,
``Maximum likelihood from incomplete data via the EM algorithm,''
\emph{Journal of the Royal Statistical Society, Series B}, 39(1):1--38, 1977.

\bibitem{Louis1982}
T. A. Louis,
``Finding the observed information matrix when using the EM algorithm,''
\emph{Journal of the Royal Statistical Society, Series B}, 44(2):226--233, 1982.

\bibitem{MengRubin1994}
X.-L. Meng and D. B. Rubin,
``On the global and componentwise rates of convergence of the EM algorithm,''
\emph{Linear Algebra and its Applications}, 199:413--425, 1994.
\end{thebibliography}

\end{document}
